%=======================================================================
\chapter{Validation statistique}

\section{Coefficient de Cohen's Kappa}

L'exactitude (\emph{Accuracy}) constitue une mesure couramment utilisée pour évaluer les
performances d'un classificateur. Toutefois, cette métrique ne tient pas compte des accords
pouvant être obtenus par hasard. Afin de fournir une évaluation plus robuste, le coefficient de
Cohen's Kappa ($\kappa$) a été calculé pour les meilleures configurations du benchmark.

Le coefficient Kappa mesure le niveau d'accord entre les prédictions du modèle et les
étiquettes réelles après correction de l'accord aléatoire~:
\begin{equation}
  \kappa = \frac{P_o - P_e}{1 - P_e}
\end{equation}
où $P_o$ représente l'accord observé et $P_e$ l'accord attendu par hasard.

Les valeurs de $\kappa$ sont généralement interprétées comme suit~:

\begin{table}[htbp]
  \centering
  \caption{Interprétation du coefficient $\kappa$.}
  \label{tab:kappa-interpretation}
  \begin{tabular}{@{}ll@{}}
    \toprule
    \thd{Valeur de $\kappa$} & \thd{Interprétation} \\
    \midrule
    $< 0{,}20$        & Accord faible \\
    $0{,}21$ -- $0{,}40$ & Accord acceptable \\
    $0{,}41$ -- $0{,}60$ & Accord modéré \\
    $0{,}61$ -- $0{,}80$ & Accord substantiel \\
    $> 0{,}80$        & Accord presque parfait \\
    \bottomrule
  \end{tabular}
\end{table}

Les résultats obtenus montrent des valeurs de Kappa élevées pour les meilleures
configurations. Le modèle \textbf{SVM avec l'ensemble Full(opt)} atteint
$\kappa = \kappaVal{}$, tandis que~:
\begin{itemize}[leftmargin=1.4em,itemsep=2pt]
  \item KNN + Full(opt)~: $\kappa = 0{,}85$~;
  \item LightGBM + Full(opt)~: $\kappa = 0{,}84$.
\end{itemize}

Ces valeurs indiquent un niveau d'accord presque parfait entre les prédictions et les classes
réelles, confirmant la qualité des modèles développés.

% -----------------------------------------------------------------------
\section{Intervalles de confiance Bootstrap}

Afin d'évaluer la stabilité des performances observées sur le jeu de test, une procédure
Bootstrap a été appliquée ($N = 2000$ rééchantillonnages avec remise).

Cette méthode permet d'estimer un intervalle de confiance à 95\,\% à partir de la distribution
des métriques recalculées sur chaque échantillon bootstrap.

Pour le meilleur modèle (\textbf{SVM + Full(opt)}), les résultats obtenus sont~:

\begin{table}[htbp]
  \centering
  \caption{Intervalles de confiance Bootstrap — SVM + Full(opt).}
  \label{tab:bootstrap-svm}
  \begin{tabular}{@{}lccc@{}}
    \toprule
    \thd{Métrique} & \thd{Valeur test} & \thd{IC 95\,\% bas} & \thd{IC 95\,\% haut} \\
    \midrule
    Accuracy & 0,920 & 0,890 & 0,950 \\
    F1-Macro & 0,9204 & 0,8895 & 0,9504 \\
    \bottomrule
  \end{tabular}
\end{table}

Ces intervalles relativement étroits témoignent d'une faible variabilité des performances et
confirment la robustesse du modèle.

Le Tableau~\ref{tab:stats-top10} synthétise les dix meilleures configurations avec Accuracy,
F1-Macro, $\kappa$ et bornes d'intervalles de confiance Bootstrap (source~:
\texttt{statistical\_validation\_summary.csv}).

\begin{table}[htbp]
  \centering
  \caption{Top 10 — Accuracy, F1-Macro, $\kappa$ et IC 95\,\% Bootstrap.}
  \label{tab:stats-top10}
  \small
  \setlength{\tabcolsep}{3pt}
  \begin{tabular}{@{}llrrrrr@{}}
    \toprule
    \thd{Ensemble} & \thd{Modèle} & \thd{Acc.} & \thd{F1} & \thd{$\kappa$} &
    \thd{F1 IC bas} & \thd{F1 IC haut} \\
    \midrule
    Full(opt) & SVM         & 0,920 & 0,920 & 0,880 & 0,890 & 0,950 \\
    D (HOG)   & SVM         & 0,903 & 0,903 & 0,855 & 0,868 & 0,935 \\
    Full(opt) & KNN         & 0,900 & 0,900 & 0,850 & 0,866 & 0,932 \\
    D (HOG)   & KNN         & 0,897 & 0,897 & 0,845 & 0,863 & 0,928 \\
    Full(opt) & LightGBM    & 0,893 & 0,894 & 0,840 & 0,859 & 0,927 \\
    Full(opt) & MLP         & 0,883 & 0,884 & 0,825 & 0,846 & 0,919 \\
    Full(opt) & XGBoost     & 0,880 & 0,880 & 0,820 & 0,843 & 0,915 \\
    D (HOG)   & MLP         & 0,877 & 0,877 & 0,815 & 0,838 & 0,912 \\
    Full(opt) & Extra Trees & 0,877 & 0,877 & 0,815 & 0,840 & 0,912 \\
    Full(opt) & LR          & 0,870 & 0,870 & 0,805 & 0,831 & 0,905 \\
    \bottomrule
  \end{tabular}
\end{table}

\begin{figure}[htbp]
  \centering
  \includegraphics[width=0.92\textwidth]{bootstrap_ci_best_fs}
  \caption{Intervalles de confiance Bootstrap (F1-Macro et Accuracy) pour les meilleurs
    modèles sur Full(opt).}
  \label{fig:bootstrap-ci}
\end{figure}

De manière générale, les meilleurs modèles présentent des intervalles de confiance fortement
chevauchants mais centrés sur des valeurs élevées, ce qui renforce la fiabilité des résultats
obtenus.

% -----------------------------------------------------------------------
\section{Analyse par validation croisée}

La validation croisée stratifiée à cinq plis utilisée durant le benchmark permet d'étudier la
stabilité des performances selon les différentes partitions des données d'apprentissage.

\subsection{Moyenne des performances}

Sur l'ensemble de \textbf{test} (300 images), le modèle SVM + Full(opt) atteint une accuracy
de \textbf{92,0\,\%} et un F1-Macro de \textbf{92,04\,\%}. Les scores moyens obtenus sur les
cinq plis de validation croisée (jeu d'apprentissage) sont légèrement inférieurs, ce qui est
attendu~: F1-Macro moyen CV = 89,5\,\%, accuracy moyenne CV = 89,5\,\%.

Ces résultats confirment que les performances observées sur le jeu de test ne sont pas dues à
une partition particulière des données.

\subsection{Écart-type des performances}

L'écart-type des scores obtenus sur les différents plis constitue un indicateur de stabilité.
Pour SVM + Full(opt)~:
\begin{itemize}[leftmargin=1.4em,itemsep=2pt]
  \item Accuracy Std (CV) = 0,0099~;
  \item F1 Std (CV) = 0,0100.
\end{itemize}

Ces valeurs particulièrement faibles indiquent que les performances restent constantes quelle
que soit la partition utilisée. Les autres modèles du Top~5 sur Full(opt) présentent
également des écarts-types inférieurs à 0,02.

Le Tableau~\ref{tab:cv-stats} résume les moyennes et écarts-types CV pour les cinq meilleurs
modèles sur Full(opt) (source~: \texttt{cv\_fold\_statistics.csv}).

\begin{table}[htbp]
  \centering
  \caption{Validation croisée (5 plis) — moyenne $\pm$ écart-type sur Full(opt).}
  \label{tab:cv-stats}
  \begin{tabular}{@{}lcc@{}}
    \toprule
    \thd{Modèle} & \thd{F1-Macro (moy. $\pm$ std)} & \thd{Accuracy (moy. $\pm$ std)} \\
    \midrule
    SVM         & $0{,}895 \pm 0{,}010$ & $0{,}895 \pm 0{,}010$ \\
    LightGBM    & $0{,}878 \pm 0{,}009$ & $0{,}878 \pm 0{,}008$ \\
    Extra Trees & $0{,}878 \pm 0{,}004$ & $0{,}879 \pm 0{,}004$ \\
    KNN         & $0{,}870 \pm 0{,}020$ & $0{,}871 \pm 0{,}020$ \\
    MLP         & $0{,}873 \pm 0{,}011$ & $0{,}873 \pm 0{,}011$ \\
    \bottomrule
  \end{tabular}
\end{table}

% -----------------------------------------------------------------------
\section{Test de McNemar}

Le test de McNemar est un test statistique non paramétrique utilisé pour comparer directement
deux classificateurs sur les mêmes observations (ici, l'ensemble de test, configuration
\textbf{Full(opt)}).

L'hypothèse nulle $H_0$ postule que les deux modèles ont des performances équivalentes. Une
valeur $p < 0{,}05$ conduit au rejet de $H_0$.

Le Tableau~\ref{tab:mcnemar} présente les comparaisons clés entre les modèles du Top~5
(source~: \texttt{mcnemar\_test\_results.csv}).

\begin{table}[htbp]
  \centering
  \caption{Test de McNemar — paires sélectionnées (Full(opt), $n = 300$).}
  \label{tab:mcnemar}
  \begin{tabular}{@{}llrrl@{}}
    \toprule
    \thd{Modèle A} & \thd{Modèle B} & \thd{$n_{01}$} & \thd{$n_{10}$} & \thd{$p$-value} \\
    \midrule
    SVM & KNN      & 20 & 14 & 0,391 \\
    SVM & LightGBM & 14 &  6 & 0,118 \\
    SVM & MLP      & 15 &  4 & 0,022 \\
    SVM & XGBoost  & 17 &  5 & 0,019 \\
    SVM & Extra Trees & 18 & 5 & 0,012 \\
    SVM & LR       & 19 &  4 & 0,004 \\
    KNN & LightGBM & 20 & 18 & 0,871 \\
    \bottomrule
  \end{tabular}
\end{table}

Les résultats montrent que SVM est significativement meilleur que LR, MLP, XGBoost et Extra
Trees ($p < 0{,}05$), mais que la différence avec \textbf{KNN} ($p = 0{,}391$) et
\textbf{LightGBM} ($p = 0{,}118$) n'est pas statistiquement significative au seuil de 5\,\%.

% -----------------------------------------------------------------------
\section{Test de Wilcoxon}

Le test des rangs signés de Wilcoxon permet de comparer les performances de deux modèles à
partir des scores F1-Macro obtenus sur les cinq plis de validation croisée.

L'hypothèse nulle est~: $H_0$~: la médiane des différences entre les deux modèles est nulle.

Pour Full(opt), les comparaisons entre SVM et les autres modèles du Top~5 donnent des
$p$-values $\geq 0{,}0625$ (seuil non atteint à 5\,\%), indiquant qu'au regard des scores CV,
les écarts observés ne sont pas statistiquement significatifs — cohérent avec le test de
McNemar pour SVM vs KNN/LightGBM.

% -----------------------------------------------------------------------
\section{Tests de Friedman et Nemenyi}

Lorsque plusieurs modèles doivent être comparés simultanément, une analyse globale a été
réalisée à l'aide du test de Friedman sur les huit ensembles de caractéristiques (A, B, C, D,
A+B, B+C, A+B+C, Full).

Pour \textbf{tous} les ensembles testés, le test de Friedman rejette $H_0$ au seuil de 5\,\%
($p < 0{,}01$), confirmant que les huit algorithmes ne sont pas équivalents. Les statistiques
les plus élevées sont observées pour B — Tex(opt) ($\chi^2 = 32{,}2$, $p < 10^{-4}$) et
A+B+C(opt) ($\chi^2 = 29{,}7$, $p = 1{,}06 \times 10^{-4}$).

En cas de rejet, un test post-hoc de Nemenyi a été appliqué pour identifier les paires de
modèles significativement différentes au sein de chaque ensemble de caractéristiques.

\begin{figure}[htbp]
  \centering
  \includegraphics[width=0.88\textwidth]{friedman_pvalues}
  \caption{Valeurs $p$ du test de Friedman par ensemble de caractéristiques (comparaison
    simultanée des 8 modèles).}
  \label{fig:friedman}
\end{figure}

\begin{table}[htbp]
  \centering
  \caption{Résultats du test de Friedman par ensemble de caractéristiques.}
  \label{tab:friedman}
  \small
  \begin{tabular}{@{}lrrl@{}}
    \toprule
    \thd{Ensemble} & \thd{Statistique} & \thd{$p$-value} & \thd{Significatif} \\
    \midrule
    B — Tex(opt)  & 32,20 & $3{,}7 \times 10^{-5}$ & Oui \\
    A+B+C(opt)    & 29,73 & $1{,}1 \times 10^{-4}$ & Oui \\
    A+B           & 28,20 & $2{,}0 \times 10^{-4}$ & Oui \\
    B+C(opt)      & 25,60 & $5{,}9 \times 10^{-4}$ & Oui \\
    A — Stats     & 24,60 & $8{,}9 \times 10^{-4}$ & Oui \\
    C — DWT(opt)  & 24,27 & $1{,}0 \times 10^{-3}$ & Oui \\
    Full(opt)     & 23,60 & $1{,}3 \times 10^{-3}$ & Oui \\
    D — HOG       & 22,33 & $2{,}2 \times 10^{-3}$ & Oui \\
    \bottomrule
  \end{tabular}
\end{table}

% -----------------------------------------------------------------------
\section{Synthèse statistique}

Les analyses statistiques réalisées dans ce projet permettent de compléter l'évaluation
classique basée sur les métriques de performance.

Les principaux enseignements sont les suivants~:
\begin{itemize}[leftmargin=1.4em,itemsep=2pt]
  \item les meilleurs modèles présentent des coefficients de Cohen's Kappa supérieurs à
        0,80, indiquant un accord presque parfait avec les annotations réelles~;
  \item les intervalles de confiance Bootstrap demeurent relativement étroits, témoignant
        d'une bonne stabilité des performances sur le jeu de test~;
  \item les faibles écarts-types observés lors de la validation croisée ($\leq 0{,}02$ pour
        le Top~5 sur Full) montrent une excellente capacité de généralisation~;
  \item le test de McNemar confirme la supériorité de SVM sur plusieurs modèles, mais pas
        sur KNN et LightGBM au seuil de 5\,\%~;
  \item le test de Friedman rejette l'équivalence des huit algorithmes sur chaque ensemble
        de caractéristiques~;
  \item les ensembles fusionnés obtiennent systématiquement des performances supérieures aux
        ensembles individuels.
\end{itemize}

Parmi l'ensemble des configurations étudiées, \textbf{SVM associé à Full(opt)} apparaît comme
la solution la plus performante et la plus robuste.

Avec une accuracy de \textbf{92,0\,\%}, un F1-Macro de \textbf{92,04\,\%}, un coefficient
$\kappa = 0{,}88$, des intervalles Bootstrap étroits et des écarts-types CV inférieurs à 0,01,
cette configuration constitue le meilleur compromis entre performance, stabilité et capacité de
généralisation.

Ces résultats valident l'efficacité de la stratégie proposée reposant sur l'optimisation des
descripteurs, la fusion multi-caractéristiques et l'évaluation comparative de plusieurs
algorithmes de Machine Learning.
